\DailyTitle{7731 Log (April 16, 2012)}

\DailySection{Progress log}

Treating leading photon as invisible, and apply the same selection as the SUSY analysis.
In addition, require all jets at least $\Delta R > 0.3$ away from the leading photon.
We can look at \MR and \R distributions.  Fortunately, data matches G+jet sample pretty
well in all cases.  Interestingly, from MC there is non-negligible contribution from QCD
sample, but data matches well with only G+jet sample.  All other SM background is
negligible in this treatment.

From MCDB (ID: 3358, 3356, 3306, 3360), the QCD samples are generated with processes
\begin{enumerate}
\item p p \textgreater ~j j
\item p p \textgreater ~j j j
\item p p \textgreater ~j j j j
\end{enumerate}
while in madgraph 1.4.3, "j" = "g u c d s u\tweakedtilde~c\tweakedtilde~d\tweakedtilde
~s\tweakedtilde", no photon.  So in principle there shouldn't
be double-counting between these two samples.  What is going on?  \ActionItem

The photon \PT spectrum matches well too.  In the inclusive selection I observe that data
do not become fully efficient until roughly 200 GeV.  But when we require that $M_R > 400,
R > 0.5$, an edge at 200 GeV of photon \PT is formed, see \ref{Figure_7731_PhotonPT}.
As a result it doesn't matter that much where we place the photon \PT cut.
There is a minimum \PT requirement of 40 GeV in the triggers, so we're lucky in this sense.
Let's put a cut at 100 GeV then.

\begin{figure}
\centering
\includegraphics[width=0.45\textwidth]{ToBeDrawn.png}
\caption{Photon \PT spectrum after minimum \MR and \R cut in the photon-invisible treatment.}
\label{Figure_7731_PhotonPT}
\end{figure}

The triggers in 2011 in PhotonHad dataset are as follows:
\begin{enumerate}
\item HLT\_Photon40\_CaloIDL\_R014\_MR150 (prescaled)
\item HLT\_Photon40\_CaloIDL\_R017\_MR500
\item HLT\_Photon40\_CaloIDL\_R023\_MR350
\item HLT\_Photon40\_CaloIDL\_R029\_MR250
\item HLT\_Photon40\_CaloIDL\_R042\_MR200
\item HLT\_Photon55\_CaloIDL\_R017\_MR500
\item HLT\_Photon55\_CaloIDL\_R023\_MR350
\item HLT\_Photon55\_CaloIDL\_R029\_MR250
\item HLT\_Photon55\_CaloIDL\_R042\_MR200
\end{enumerate}
Last four triggers are backup triggers it seems.  With the treatment of photon-invisible,
\MR will decrease, while \R will increase.  Since we're requiring photons at least 100 GeV,
in terms of \MR threshold probably all triggers are fine.  The relevant question is then how
much \R increases.  And it increased by a lot.  This thing couldn't possibly have worked.
What is happening?

The other triggers are listed here:
\begin{enumerate}
\item HLT\_Photon60\_CaloIdL\_HT300 (prescaled)
\item HLT\_Photon60\_CaloIdL\_MHT70 (prescaled)
\item HLT\_Photon60\_CaloIdXL\_HT400
\item HLT\_Photon60\_CaloIdXL\_HT500
\item HLT\_Photon60\_CaloIdXL\_MHT100
\item HLT\_Photon60\_CaloIdXL\_MHT90
\item HLT\_Photon90EBOnly\_CaloIdVL\_IsoL\_TriPFJet25
\item HLT\_Photon90EBOnly\_CaloIdVL\_IsoL\_TriPFJet30
\end{enumerate}
Maybe it's one of these triggers... \ActionItem

Anyways, in the plot with $R^2 > 0.20$, \MR seems to be fully efficient above 300-ish, and
in the plot with $M_R > 400$, \RSquare seems to be fully efficient above 0.15-ish.  So in our
region of interest it seems to be fine.





